\section{Estado del arte y trabajos relacionados}

La consulta de información técnica puede abordarse mediante distintas estrategias. La opción más sencilla consiste en revisar manualmente los documentos. También pueden utilizarse buscadores por palabras clave, bases de datos relacionales, chatbots conectados a un modelo de lenguaje o sistemas RAG. La diferencia principal se encuentra en la forma en que cada alternativa interpreta la pregunta, recupera la información y permite verificar la respuesta.

En una organización dedicada a la innovación, la solución debe considerar que el corpus puede crecer y cambiar. Un mecanismo útil no debe limitarse a encontrar una coincidencia de texto; debe conservar la relación entre la respuesta y el documento de origen. También debe permitir que un responsable detecte cuándo la pregunta no puede responderse con la información disponible.

\subsection{Búsqueda manual de procedimientos}

La búsqueda manual consiste en explorar carpetas, archivos y páginas hasta encontrar la instrucción requerida. Su ventaja es que la persona revisa directamente el documento y puede interpretar elementos que un sistema automático no reconoce. Además, no requiere infraestructura adicional.

Su principal desventaja es el tiempo. El usuario debe conocer la ubicación, el nombre o la estructura del archivo. Cuando existen documentos similares o versiones antiguas, aumenta la posibilidad de consultar material incorrecto. El resultado también depende de la experiencia de cada persona, por lo que el proceso no siempre es uniforme.

\subsection{Búsqueda por palabras clave}

La búsqueda por palabras clave permite localizar coincidencias dentro de varios documentos. Es fácil de implementar, tiene un costo bajo y funciona bien cuando la colección es pequeña y el usuario conoce los términos exactos. Puede complementarse con filtros por fecha, tipo de archivo o área.

La limitación es que la coincidencia literal no comprende necesariamente el significado. Una pregunta con sinónimos, abreviaturas o una estructura distinta puede no devolver el documento correcto. También puede producir demasiados resultados si una palabra aparece en muchas secciones, sin explicar cuál de ellos responde realmente a la pregunta.

\subsection{Chatbot sin recuperación documental}

Un chatbot conectado directamente a un modelo de lenguaje ofrece una interacción natural. El usuario puede preguntar sin conocer la organización de los documentos y recibe una respuesta redactada con rapidez. Esta alternativa es útil para tareas generales de explicación o generación de texto.

Para procedimientos técnicos, el chatbot sin recuperación presenta riesgos. Puede usar conocimiento externo al dominio, mezclar versiones, inventar pasos o no ser capaz de señalar la fuente. La respuesta puede ser lingüísticamente correcta y, al mismo tiempo, no corresponder al procedimiento oficial. Por esta razón, se considera una solución incompleta para el problema de este proyecto.

\subsection{Sistemas de recuperación aumentada por generación}

Los sistemas RAG agregan una etapa que busca fragmentos antes de consultar al modelo. El buscador semántico puede identificar información relacionada aunque la pregunta y el documento no utilicen exactamente las mismas palabras. El modelo recibe los fragmentos como contexto y puede generar una explicación con referencias.

Esta alternativa introduce nuevos retos. La calidad de la respuesta depende de la ingesta, la segmentación, los embeddings, el motor vectorial, el valor de \texttt{top\_k}, el prompt y la evaluación. Si el recuperador no encuentra el fragmento adecuado, el modelo puede recibir un contexto incompleto. Por ello, el sistema debe conservar las puntuaciones y establecer reglas para declarar insuficiencia.

\subsection{Copilotos documentales en entornos técnicos}

Un copiloto documental se diferencia de un chatbot general porque trabaja con un dominio, un corpus y una responsabilidad definida. Puede responder preguntas frecuentes, resumir instrucciones y dirigir al usuario hacia la fuente. La decisión final permanece en la persona responsable, especialmente cuando la información afecta seguridad, calidad, operación o cumplimiento.

El diseño propuesto para CII.IA añade una separación por componentes. La interfaz se concentra en la experiencia de consulta, el backend controla el acceso al proveedor y el agente coordina el flujo. El corpus y los metadatos son una capa independiente que puede cambiar sin obligar a rehacer la interfaz.

\subsection{Criterios de comparación}

Para comparar las alternativas se consideran cinco criterios. El primero es la facilidad de consulta, es decir, si el usuario puede preguntar con lenguaje natural. El segundo es la comprensión semántica. El tercero es la trazabilidad de la información. El cuarto es el esfuerzo de implementación y mantenimiento. El quinto es la capacidad de controlar respuestas cuando no existe información suficiente.

\subsection{Tabla comparativa}

La Tabla \ref{tab:comparativa} presenta la comparación de las opciones identificadas.

\renewcommand{\arraystretch}{1.35}
\begin{longtable}{p{3.3cm}p{3.4cm}p{3.4cm}p{3.5cm}}
\caption{Comparación de alternativas para consultar documentación técnica.}\label{tab:comparativa}\\\toprule
\textbf{Alternativa} & \textbf{Ventajas} & \textbf{Limitaciones} & \textbf{Fuentes y trazabilidad}\\\midrule
\endfirsthead
\toprule\textbf{Alternativa} & \textbf{Ventajas} & \textbf{Limitaciones} & \textbf{Fuentes y trazabilidad}\\\midrule
\endhead
\bottomrule\endfoot
Búsqueda manual & No requiere desarrollo y permite revisar el documento completo. & Consume tiempo y depende de conocer dónde buscar. & La fuente es visible, pero el proceso no es uniforme.\\
Búsqueda por palabras clave & Es sencilla y rápida en documentos pequeños. & No comprende sinónimos, contexto ni lenguaje natural. & Puede mostrar coincidencias sin una explicación.\\
Chatbot sin RAG & Ofrece interacción natural y respuestas redactadas. & Puede inventar información o utilizar conocimiento externo al corpus. & La fuente no queda garantizada.\\
Copiloto RAG propuesto & Combina pregunta natural, contexto documental y respuesta resumida. & Depende de la calidad de la ingesta, recuperación y modelo. & Puede devolver archivo, página, sección, fragmento y puntuación.\\
Atención humana & Permite interpretar casos complejos y validar decisiones. & Tiene disponibilidad limitada y no siempre responde de inmediato. & Depende del registro de la consulta y del responsable.\\
\end{longtable}

La comparación muestra que la alternativa RAG ofrece el mejor equilibrio para el objetivo del proyecto, siempre que se complemente con una fuente documental preparada y mecanismos de control. No elimina la revisión humana, pero puede disminuir el tiempo de localización y proporcionar una primera respuesta sustentada.

El resultado de esta comparación orienta las decisiones del sistema propuesto. La interfaz debe ser sencilla como un chatbot, el backend debe conservar la seguridad de una API propia y el agente debe utilizar la trazabilidad de una base documental. El desarrollo posterior debe verificar si estas ventajas se mantienen con preguntas reales y con el corpus que será entregado por el Integrante 1.
